\section{Theory}\label{sec:theory}

\subsection{Common setup}

\noindent Let $A$ be an individual's additive breeding value for a transferable capacity and $C=A+E$ its phenotype, with
$A\sim\N(\mu,V_A)$ and $E\sim\N(0,V_E)$ independent, $V_P=V_A+V_E$, and $h^2=V_A/V_P$. Observers see an averaged score from
$n$ events,
\begin{equation}
Y=\lambda C+U+\sigma\sqrt{\rho}\,Z_0+\sigma\sqrt{(1-\rho)/n}\,Z_n,\label{eq:score}
\end{equation}
where $U\sim\N(0,\tau^2)$ is persistent nuisance (training, opportunity), $\rho$ is the shared fraction of event error, and
$Z_0,Z_n$ are independent standard normals. Target validity is
\begin{equation}
R=\Corr(C,Y)^2=\frac{\lambda^2V_P}{\lambda^2V_P+\tau^2+\sigma^2[\rho+(1-\rho)/n]}.\label{eq:validity}
\end{equation}
Let $Z=(Y-\E Y)/\SD(Y)$ denote \emph{standing}, the individual's standardized position in the assessed ranking. Transmission
is Mendelian and unbiased within families, and inheritance follows the infinitesimal model~\cite{bulmer1980}. Response is the
change in mean breeding value among offspring, computed with the Price equation~\cite{price1970}.

\subsection{The core identity: information times opportunity}

Let reproductive allocation within a sex depend only on the observed score, $w=W(Y)\ge0$ with $\E W>0$.

\begin{proposition}[Information--opportunity identity]\label{prop:identity}
If\/ $\E[A\mid Y]$ is linear in $Y$ (exact under joint Gaussianity), the selection differential on breeding value is
\begin{equation}
S=\frac{\Cov(A,W)}{\E W}=\underbrace{\sigma_A}_{\text{genetic}}\;\underbrace{h\sqrt{R}}_{\text{accuracy}}\;
\underbrace{\Corr(Y,W)}_{\text{alignment}}\;\underbrace{\CV(W)}_{\text{opportunity}},\label{eq:identity}
\end{equation}
where $h\sqrt{R}=\Corr(A,Y)$ is the accuracy of the score as a predictor of breeding value and $\CV(W)^2$ is the opportunity
for selection attributable to standing~\cite{crow1958}. Without linearity, $|S|\le\SD(\E[A\mid Y])\,\CV(W)$.
\end{proposition}

Equation~\eqref{eq:identity} is the animal breeder's key equation, gain equals intensity times accuracy times
$\sigma_A$~\cite{lush1937,hazel1943,falconer1996}, with intensity $i=\Corr(Y,W)\CV(W)$; the bound is the opportunity bound of
sexual-selection theory~\cite{crow1958,moorad2013,henshaw2016}. None of this is new mathematics. Its value is organizational:
each ingredient of the hypothesis maps onto one factor. The performance ecology sets accuracy through $R$; institutions
(mating system, marriage rules, how strongly standing governs access) set intensity; and assortment changes $\sigma_A$ across
generations. Three consequences follow.

\emph{Information and inequality are separable.} $R$ enters only through accuracy. Any allocation that depends on
standardized standing alone has $\CV(W)$ and $\Corr(Y,W)$ independent of $R$, so the response scales exactly as $\sqrt{R}$ at
fixed opportunity. A contrast that raises reproductive inequality is an opportunity effect, not an information effect, and
must not be counted as one.

\emph{Equal contribution is a hard null.} If every participating parent has the same expected contribution, $\CV(W)=0$ and
$S=0$ for any $R$. Information can change which individuals reproduce only if contribution varies.

\emph{Leverage is greatest where traits are hardest to see.} The proportional gain from raising validity from $R_0$ to $R_1$
is $\sqrt{R_1/R_0}$, largest when baseline validity is low. In U.S.\ Holstein cattle, replacing progeny testing, which
waits for daughters' lactation records, with genomic prediction at birth cut generation intervals sharply and raised annual
genetic gain by 50--100\% for yield and three- to fourfold for lowly heritable traits~\cite{garciaruiz2016}. One difference
matters: in hominins earlier information did not shorten generations, because life histories
slowed~\cite{kaplan2000,zollikofer2024}. The leverage is on accuracy at the moment decisions are made, not on generation
time. This is the precise sense in which performance comes ``before outcomes'' (claim N4 of Table~\ref{tab:audit}).

\begin{remark}[Phenotypic validity is not genetic accuracy]
The accuracy that enters the response is $\Corr(A,Y)$. It equals $h\sqrt R$ only when the non-genetic part of the score is
independent of breeding value. Where it is not, as when access to training is inherited with rank, a score can predict the
phenotype better and the breeding value worse; selection-index theory gives the general case~\cite{hazel1943}. Every statement
about validity in this paper concerns the phenotype. Genetic accuracy has to be measured and cannot be assumed from it.
\end{remark}

\subsection{Model I: three routes from information to selection}

Two routes do not fit the form of Eq.~\eqref{eq:identity}: fitness can depend on one's own capacity directly, and on a
partner's capacity. Individuals of sex $s$ pair by ideal reciprocal rank matching on $Y$ (or do not pair, in which case the
partner route is absent). Expected reproductive contribution is
\begin{equation}
w_s=\exp(\beta_{N,s}C+\beta_{D,s}C_p+\beta_{S,s}Z),\label{eq:routes}
\end{equation}
where $C_p$ is the partner's capacity. The three coefficients are distinct routes: $\beta_N$ is the return to one's own
capacity (provisioning, survival of offspring); $\beta_D$ is the return to a partner's capacity (direct benefits of a capable
partner); and $\beta_S$ is the return to one's own standing (earlier marriage, more or preferred mates, kin and ally
support, prestige).

\begin{proposition}[Three-route decomposition]\label{prop:routes}
Under Gaussianity and ideal rank matching on $Y$, the within-sex selection differential on breeding value is
\begin{equation}
S_s=V_A\Big[\beta_{N,s}+\beta_{D,s}R+\beta_{S,s}\sqrt{R/V_P}\,\Big],\label{eq:threeroute}
\end{equation}
and the response of mean breeding value is $\Delta\bar A=\tfrac12(S_f+S_m)$.
\end{proposition}

\begin{corollary}[Information bottleneck]\label{cor:bottleneck}
If $R=0$, only the own-capacity route remains. However strong the preference for capacity, choice generates no selection on
capacity without information about it.
\end{corollary}

\begin{corollary}[Selection follows the score]\label{cor:score}
Suppose the score also contains a heritable display component $D$ with additive variance $V_{A,D}$, independent of $C$.
Standing selection changes the display by $\beta_SV_{A,D}/\SD(Y)$ and capacity by $\beta_S\lambda V_A/\SD(Y)$. When display
variance dominates the score, choice drives presentation rather than capacity.
\end{corollary}

Three points deserve emphasis. First, the standing route is a pure information effect: $\Var(\ln w)=\beta_S^2$ whatever $R$
is. The partner-quality route is not. Assortment raises $\Var(X)=V_P(1+R)/2$ and therefore the inequality of pair output;
holding $\Var(\ln W)$ fixed in a model where selection acts only on pair output, the gain from $R=0$ to $R=1$ is $\sqrt2$,
not 2: on a log scale, half of the apparent information effect is increased reproductive inequality. Second, the standing route scales with
$\sqrt{R}$ while the partner-quality route scales with $R$, so standing is the less information-limited channel: at $R=0.25$
(with $V_P=1$) the partner route operates at a quarter of its full-information strength and the standing route at half. Third,
the standing route does not require assortment at all. Status-dependent mating, earlier pairing of higher-ranked
individuals (the Darwin--Fisher mechanism~\cite{kirkpatrick1990}, observed among Tsimane~\cite{vonrueden2011}), and kin or ally
support all enter through $\beta_S$.

\subsection{Model II: portfolio assessment}\label{sec:portfolio}

Let performance in domain $k$ be $T_k=G+S_k$, where $G\sim\N(0,V_G)$ is a general capacity and $S_k\sim\N(0,V_S)$ is
domain-specific, all independent. Observers see $n$ events in each of $K$ domains with domain-persistent nuisance $U_k$
(variance $\tau^2$), within-domain shared error fraction $\rho$, and a general nuisance $\Omega$ (variance $\omega^2$) common
to all domains, such as opportunity or presentation:
\begin{equation}
Y=\frac1K\sum_{k=1}^K\big(T_k+U_k+\bar e_k\big)+\Omega.\label{eq:portfolio}
\end{equation}

\begin{proposition}[Portfolio validity]\label{prop:portfolio}
Validity for the general capacity is
\begin{equation}
R_G=\frac{V_G}{V_G+\big(V_S+\tau^2+\sigma^2[\rho+(1-\rho)/n]\big)/K+\omega^2}.\label{eq:RG}
\end{equation}
Repetition within one domain has ceiling $V_G/(V_G+V_S+\tau^2+\rho\sigma^2+\omega^2)$; diversity raises the ceiling to
$V_G/(V_G+\omega^2)$.
\end{proposition}

\begin{proposition}[Diversity concentrates selection on general capacity]\label{prop:diversity}
Under standing selection $w=\exp(\beta_SZ)$, with heritabilities $h_G^2,h_S^2,h_\Omega^2$, the selection differentials on
breeding values are
\begin{equation}
S_{A_G}=\frac{\beta_Sh_G^2V_G}{\SD(Y)},\quad S_{A_{S_k}}=\frac{\beta_Sh_S^2V_S}{K\,\SD(Y)},\quad
S_{A_\Omega}=\frac{\beta_Sh_\Omega^2\omega^2}{\SD(Y)}.\label{eq:diversity}
\end{equation}
\end{proposition}

For a fixed observation budget, spreading events across domains lowers $\SD(Y)$ and so raises selection on the general
capacity, while the push on each specialized skill falls as $1/K$. Diversity cannot remove the general nuisance $\Omega$,
which sets the irreducible ceiling and remains the channel for selection on presentation. The prediction is concrete:
composite indices across many performance domains should predict reproductive outcomes better than any single domain, as in
budgerigars~\cite{medinagarcia2021}. This is also the answer to objection O29. That a hunter's reputation tracks his measured
skill only weakly~\cite{stibbardhawkes2018} is what Eq.~\eqref{eq:RG} gives for $K=1$ with large event noise; the question for
the hypothesis is what assessors achieve with many venues and many years.

\begin{remark}[Reputation pooled by talk]
If each assessor also hears reports from $g$ others about independent events, the events behind a judgment rise from $n$ to
$gn$ at almost no cost to the performer, but hearsay adds a shared distortion of variance $\varsigma^2$ that no number of
informants removes: the bracket in Eq.~\eqref{eq:validity} becomes $\tau^2+\varsigma^2+\sigma^2[\rho+(1-\rho)/(gn)]$. Language
therefore drives validity quickly to a ceiling set by shared bias. That ceiling is where the rival mechanisms of copying and
prestige cascades~\cite{duval2023,gouda2018} act: they inflate $\varsigma^2$.
\end{remark}

\subsection{Model IIb: the generalist target}\label{sec:generalist}

Proposition~\ref{prop:diversity} has a consequence that deserves separate statement, because it identifies what kind of mind
this form of selection builds. Write $V_{A,G}=h_G^2V_G$ and $V_{A,S}=h_S^2V_S$.

\begin{corollary}[Generalist aim]\label{cor:generalist}
When all candidates are compared on the same $K$ domains,
\begin{equation}
\frac{S_{A_G}}{S_{A_{S_k}}}=K\,\frac{V_{A,G}}{V_{A,S}}.\label{eq:aimK}
\end{equation}
A single venue ($K=1$) cannot tell a generalist from a specialist: it selects the two kinds of ability in proportion to their
genetic variances. Each added venue dilutes every specialty and leaves the pull on general capacity undiminished.
\end{corollary}

Two qualifications keep this honest. The push summed over all specialties does not fall; it is spread over $K$ independent
directions while the push on general capacity stays in one. And the general nuisance $\Omega$ is selected in the same
proportion as $G$ (Eq.~\ref{eq:diversity}): what a portfolio rewards is whatever is common to all venues, which includes
general cognitive capacity but also health, condition and opportunity. $V_{A,G}$ should therefore be read broadly. That
breadth is itself informative. Most genetic influences on cognitive abilities are shared across abilities~\cite{kovas2006,plomin2015},
and a large mutational target makes a general factor a natural index of genome-wide condition~\cite{rowe1996,keller2006}.

\begin{proposition}[Breadth buys time]\label{prop:time}
Let $R^*<V_G/(V_G+\omega^2)$ be a target validity for general capacity, $D=V_G(1/R^*-1)-\omega^2$ and
$Q=V_S+\tau^2+\rho\sigma^2$. The number of observed events per domain needed to reach the target is
\begin{equation}
n^*(K)=\frac{(1-\rho)\sigma^2}{KD-Q}\quad\text{for }K>K_{\min}=\frac{Q}{D},\label{eq:nstarK}
\end{equation}
and the target is unreachable for $K\le K_{\min}$ however long observation continues.
\end{proposition}

With the illustrative values used in the numerical experiments and $R^*=0.5$, no amount of watching reaches the target with
five domains or fewer; six domains need \PolyNSix{} events each, eight need \PolyNEight, and twelve need about one
(Fig.~\ref{fig:polymath}A). This joins the two halves of the hypothesis. Mastery of any one demanding skill arrives late, so a
specialist's quality is known late. A generalist's can be known early, because a little evidence from each of many domains,
including how quickly new ones are picked up, isolates the one thing the domains share. Performance before outcomes is
feasible for general capacity in a way that it is not for any single skill.

\begin{proposition}[Expenditure follows breadth]\label{prop:polymath}
Let a heritable general investment $g$ raise performance in every domain and heritable specific investments $s_k$ raise it in
one, $T_k=G+g+S_k+s_k$, at costs $\tfrac12c_Gg^2$ and $\tfrac12c_Ss_k^2$ in log fitness, under standing selection
$w\propto\exp(\beta_SZ)$. The evolutionary equilibrium is
\begin{equation}
g^*=\frac{\beta_S}{c_G\,\SD(Y)},\qquad s_k^*=\frac{\beta_S}{K\,c_S\,\SD(Y)},\label{eq:gstar}
\end{equation}
and the share of total investment cost devoted to general capacity is
\begin{equation}
\pi_G=\frac{Kc_S}{Kc_S+c_G},\qquad\text{which passes one half at }K^\dagger=\frac{c_G}{c_S}.\label{eq:share}
\end{equation}
\end{proposition}

General capacity is the expensive kind of ability if it resides in tissue that serves every domain, so $c_G/c_S$ is large.
Proposition~\ref{prop:polymath} then says that general capacity dominates investment only once the number of commonly assessed
domains exceeds that cost ratio. A lineage with a narrow repertoire should invest in specialized skill however intense its
competition; one with a broad repertoire should invest in general capacity. As cumulative culture adds domains (tools, fire,
song, speech, craft, healing), $K$ rises and with it the weight on $g$, which ties the demand for a general-purpose brain to
the breadth of culture and not to its depth. As history this is conjecture. As a comparison across species and societies it is
testable (prediction P22).

\begin{proposition}[Division of labor restores the specialist]\label{prop:division}
Suppose each individual performs, and is compared with others, only in the domain of their greatest specific talent. Within a
specialty the covariances of the score with breeding values satisfy
\begin{equation}
\frac{\Cov(A_G,Y)}{\Cov(A_{S},Y)}=\frac{V_{A,G}}{c_K\,V_{A,S}},\label{eq:division}
\end{equation}
where $c_K$ is the variance of the largest of $K$ independent standard normal variables ($c_2=0.68$, $c_5=0.45$, $c_{10}=0.34$,
$c_{20}=0.28$). Relative to common comparison (Eq.~\ref{eq:aimK}) the aim on general capacity is weaker by the factor $Kc_K$:
\PolyDivFive-fold for five domains, \PolyDivTen{} for ten and \PolyDivTwenty{} for twenty. The instrument is also a single venue,
with the low validity ceiling of Proposition~\ref{prop:portfolio}.
\end{proposition}

\begin{remark}[Commensurability, not the rule]
Honoring each person's best performance, instead of their average, does not by itself change the aim. By Stein's lemma
$\Cov(G,\max_kY_k)=V_G$ and $\Cov(S_k,\max_jY_j)=V_S/K$, so the ratio in Eq.~\eqref{eq:aimK} is unchanged as long as everyone is
seen in every domain; only validity falls, because a maximum is a noisier instrument than a mean. What removes the generalist
aim is the loss of a common set of domains.
\end{remark}

The setting in which selection is aimed at the generalist is therefore specific: small societies with little division of labor
beyond sex and age, in which every adult forages, makes tools, builds, sings, argues and raises children under the eyes of the
same few dozen people. Societies with occupational specialization narrow each person's publicly comparable performance toward
one domain and, on this model, relax selection on general capacity in favor of whatever each niche rewards. That is one
candidate reading of smaller Holocene brains in more complex societies~\cite{desilva2021}, beside smaller bodies and inherited
rank; the reduction itself is disputed (Section~\ref{sec:present}).

\subsection{Model III: assortment across generations}

Under ideal rank matching on scores the phenotypic correlation between partners is $m=R$; with idiosyncratic tastes $m<R$
(Section~\ref{sec:experiments}). Because $A$ affects pairing only through $C$, $\Cov(A_f,A_m)=mh^2V_A$.

\begin{proposition}[Assortment inflates additive variance]\label{prop:assort}
In the infinitesimal model with constant linkage-equilibrium (genic) variance $V_{LE}$ and log-linear selection,
\begin{equation}
V_{A,t+1}=\tfrac12V_{A,t}\big(1+mh_t^2\big)+\tfrac12V_{LE}.\label{eq:assort}
\end{equation}
For $m<1$ the equilibrium $V_A^*$ is the positive root of $(1-m)V^2+(V_E-V_{LE})V-V_{LE}V_E=0$; for $m=1$ a finite
equilibrium exists only if $V_E>V_{LE}$. Under the pair-output model the long-run per-generation response is
$\tfrac b2V_A^*(1+m)$, an amplification over random mating of
\begin{equation}
F(m)=\frac{V_A^*}{V_{LE}}(1+m).\label{eq:F}
\end{equation}
\end{proposition}

Assortment builds positive gametic-phase disequilibrium among capacity-increasing alleles~\cite{bulmer1980,crow1968}.
Table~\ref{tab:F} gives $F$ over plausible ranges. At the modern partner correlation for IQ score ($m=0.44$~\cite{horwitz2023})
and genic heritability 0.4--0.5, $F\approx\MGFmod$; at $m=0.8$ it ranges from \MGFhiLo{} to \MGFhiHi. Modern correlations are
not ancestral values, but they show the scale is not negligible.

\begin{table}[t]
\caption{Long-run amplification $F(m)$ relative to random mating.}\label{tab:F}
\footnotesize
\begin{tabularx}{\columnwidth}{@{}Xccccc@{}}
\toprule
 & \multicolumn{5}{c}{Partner correlation $m$}\\
\cmidrule(l){2-6}
Genic $h^2$ & 0.2 & 0.4 & 0.6 & 0.8 & 1.0\\
\midrule
\tabinput{generated/tab_F.tex}
\bottomrule
\end{tabularx}
\tablenote{Genic $h^2=V_{LE}/(V_{LE}+V_E)$. $\infty$: no finite equilibrium in the infinitesimal model; finite loci,
stabilizing selection, or mutation--selection balance would bind first.}
\end{table}

\subsection{Model IV: positional demand, energy budget, and tissue}

Now let capacity be costly. Individual fitness is
\begin{equation}
w=\exp\!\Big(\beta_NC+\beta_SZ-\frac c2C^2\Big),\label{eq:cost}
\end{equation}
where $c$ is a quadratic cost (energy, development time, risk).

\begin{proposition}[Positional persistence]\label{prop:persist}
Under Gaussianity the exact one-generation response and the equilibrium are
\begin{align}
\Delta\bar A&=\frac{V_A\big(\beta_N+\beta_S\sqrt{R/V_P}-c\mu\big)}{1+cV_P},\label{eq:persist}\\
\mu^*&=\frac{\beta_N+\beta_S\sqrt{R/V_P}}{c}.\label{eq:mustar}
\end{align}
The standing gradient does not depend on $\mu$: as the whole population improves, the bar rises with it. The equilibrium
does not depend on $V_A$, which sets only the approach rate $V_Ac/(1+cV_P)$.
\end{proposition}

\begin{corollary}[Positional surplus]\label{cor:surplus}
Standing is zero-sum: $\E[\beta_SZ]=0$ at every $\mu$. The equilibrium therefore exceeds the level $\beta_N/c$ that maximizes
mean log fitness by $\beta_S\sqrt{R/V_P}/c$, and the population bears the cost of that surplus without aggregate benefit,
unless capacity also generates benefits for others.
\end{corollary}

The surplus is the evolutionary analogue of an arms race, and it has a constructive reading. If capacity produces spillovers
(innovations, teaching, coordination) that raise others' fitness, competition for standing becomes a private engine for a
public good, a link to cultural-brain accounts~\cite{muthukrishna2018} that deserves formal treatment.

Contrast an absolute, saturating ecological benefit $f(C)=\beta_E\kappa(1-e^{-C/\kappa})$. Its gradient
$\E[f'(C)]=\beta_Ee^{V_P/2\kappa^2}e^{-\mu/\kappa}$ declines as the population improves, and the elasticity of the equilibrium
with respect to the inverse cost is $1/(1+\mu^*/\kappa)<1$, falling as investment grows. Under positional demand the
elasticity is exactly one. Ecological demand that never saturates behaves like positional demand; the distinction is
saturation, not the label ``ecological.'' The cultural ratchet, which keeps creating new niches, is one way ecological
demand avoids saturation~\cite{muthukrishna2018}; rank-based assessment is another, and the two are complementary.

\emph{The anatomical bridge.} Capacity need not be made of tissue. Let $C=aB+K+E$, where $B$ is tissue investment (brain
size), $K$ is tissue-independent organization (efficiency, wiring, motivation, plasticity), and $a\ge0$ is the functional
contribution of tissue to capacity. Both $B$ and $K$ are heritable with additive covariance matrix $\mathbf G$, and each
carries a quadratic cost:
\begin{equation}
\ln w=\beta_{\mathrm{eff}}C-\frac{\kappa_B}2B^2-\frac{\kappa_K}2K^2,\qquad\beta_{\mathrm{eff}}=\beta_N+\beta_S\sqrt{R/V_P}.\label{eq:tissuefit}
\end{equation}

\begin{proposition}[Tissue tracks demand, information, and budget]\label{prop:tissue}
Under Lande's multivariate response $\Delta\bar{\mathbf z}=\mathbf G\bm\beta$~\cite{lande1979,lande1983} with nonsingular
$\mathbf G$, the unique equilibrium is
\begin{equation}
B^*=\frac{a\beta_{\mathrm{eff}}}{\kappa_B},\qquad K^*=\frac{\beta_{\mathrm{eff}}}{\kappa_K},\qquad\frac{B^*}{K^*}=\frac{a\kappa_K}{\kappa_B},\label{eq:tissue}
\end{equation}
independent of $\mathbf G$. If $a=0$, tissue returns to its cost optimum whatever the genetic correlation between tissue and
capacity; correlated responses are transient.
\end{proposition}

\begin{corollary}[Information and energy co-limit tissue]\label{cor:colimit}
At current tissue $\bar B$, the selection gradient on tissue is positive only if
\begin{equation}
\sqrt{R/V_P}>\frac{\kappa_B\bar B/a-\beta_N}{\beta_S},\label{eq:threshold}
\end{equation}
a validity threshold that rises with current tissue. Taking the assessed capacity to be the general factor of
Proposition~\ref{prop:portfolio}, validity cannot exceed $R_{\max}=V_G/(V_G+\omega^2)$, so maximal tissue investment is
$B_{\max}=a(\beta_N+\beta_S\sqrt{R_{\max}/V_P})/\kappa_B$: the energy budget and the general nuisance jointly cap the organ.
\end{corollary}

Proposition~\ref{prop:tissue} turns a caveat into conditions. Permanent enlargement requires that larger tissue functionally
improve the assessed capacity ($a>0$). A genetic correlation alone, such as the $r_g\approx0.24$ between brain volume and
intelligence~\cite{jansen2020}, is compatible with $a>0$ but does not establish it; the within-family association of brain
size with ability~\cite{lee2019} and the response to artificial selection in guppies~\cite{kotrschal2013} are better evidence.
Tissue investment is proportional to demand, to $\sqrt R$, and to the inverse of tissue cost. As tissue becomes cheaper
through cooking, provisioning, or cooperative care, allocation shifts toward tissue; when it becomes dearer, toward
organization, consistent with later changes in brain shape without further size increase~\cite{neubauer2018}. With the
illustrative values used below the approach time is $1/(G_{BB}\kappa_B)$, about $10^2$--$10^3$ generations, short relative to
fossil sampling intervals. Persistent positional demand combined with stepwise relaxation of the budget would therefore
produce stasis punctuated by steps. That is the kind of pattern found in Section~\ref{sec:explananda}, and it is compatible
with random-walk and stasis fits~\cite{hubbe2026}; but the model is given the dates of its steps and does not predict them,
so the match in shape is a consistency check and not a test (prediction P9 states what a test requires).
Corollary~\ref{cor:colimit} adds a second lever. With $\beta_N=0$, $\beta_S\sqrt{R/V_P}=0.04$ at $R=0.5$, $a=0.5$, and
$\kappa_B=0.0025$, the tissue cap is about 5 standard deviations when assessment is limited to one contest domain
($R_{\max}=0.19$) but about 10 with a diverse portfolio ($R_{\max}=0.83$), with no change in the energy budget. Wider and fairer
public assessment can raise the ceiling on the organ just as cheaper calories can.

\subsection{Model V: contest, coalition, and the aim of selection}\label{sec:modelv}

Popular accounts give physical contest and coalitional fighting the leading roles (Section~\ref{sec:audit}). Three
results state what such contests can and cannot do. They concern, in turn, what a contest selects, how much it reveals, and
how much reproduction it governs.

\subsubsection{Aim} Let formidability $F$ (size, strength, fighting skill) and general capacity $G$ be independent
heritable traits with phenotypic variances $V_F,V_G$ and heritabilities $h_F^2,h_G^2$. Standing rests on a composite of
venues $k$, weighted $\omega_k\ge0$:
\begin{equation}
Y=\sum_k\omega_k\big(\ell_{Fk}F+\ell_{Gk}G+e_k\big),\label{eq:venues}
\end{equation}
with independent errors $e_k$. Write $\bar\ell_F=\sum_k\omega_k\ell_{Fk}$ and $\bar\ell_G=\sum_k\omega_k\ell_{Gk}$ for the
composite loadings.

\begin{proposition}[The score sets the aim]\label{prop:aim}
Under joint Gaussianity, for any allocation $W(Y)$ with intensity $i=\Corr(Y,W)\CV(W)$,
\begin{equation}
S_{A_F}=\frac{i\,h_F^2\bar\ell_FV_F}{\SD(Y)},\qquad S_{A_G}=\frac{i\,h_G^2\bar\ell_GV_G}{\SD(Y)},\label{eq:aim}
\end{equation}
so that $S_{A_G}/S_{A_F}=\bar\ell_Gh_G^2V_G/(\bar\ell_Fh_F^2V_F)$ whatever the intensity, the shape of the allocation rule, or
the amount of noise.
\end{proposition}

Intensity scales the response; the loading aims it. A society in which standing is won by force has a large $\bar\ell_F$ and
a small $\bar\ell_G$, and however fierce the competition it selects mostly on the body. This is what the comparative record
shows: where males contest most intensely, bodies and canines enlarge~\cite{plavcan2012} and relative brain size does
not~\cite{schillaci2006,lemaitre2009,fitzpatrick2012}. Selection turns toward capacity when the weights $\omega_k$ move toward
venues that load on it: hunting returns, craft, mediation, oratory, teaching. The shift from standing taken by dominance to
standing conferred as prestige~\cite{henrich2001,cheng2013} is such a move.

\begin{corollary}[A cue to a state is not a target]\label{cor:cue}
Let a preferred cue be $Q=T+A_Q+e$, where $T$ is a transient state through which every individual passes (age, parity) and
$A_Q$ is a heritable component independent of capacity. Then $S_{A_Q}=i\,V_{A_Q}/\SD(Q)$ and $S_{A_G}=0$. The larger the share
of cue variance due to state, the smaller the heritable response; and the strength of the preference is irrelevant to
capacity.
\end{corollary}

This disposes of claims N9 and N10 as engines of brain evolution. Preferences for youth are strong and nearly
universal~\cite{kenrick1992}, but to the extent that they track age they select a stage of life, not a lineage. Heritable
variation in the \emph{timing} of life-history events is a different matter and can be selected, as it is at present
(Section~\ref{sec:gss}).

\subsubsection{Repeatability} Animal contests are sequences of mutual assessment in which each exchange is a
sample~\cite{enquist1983,arnott2009}. The same holds for an audience. Let each contest event cost its participants $\delta>0$
in expected log fitness (injury, death, time), and let the value of assessment to those who use it be proportional to its
accuracy, $B\sqrt{R(n)}$, as it is when the expected quality of a chosen partner or ally is intensity times accuracy.

\begin{proposition}[Safe contests inform more]\label{prop:repeat}
With $R(n)$ from Eq.~\eqref{eq:validity}, the net benefit $B\sqrt{R(n)}-\delta n$ has a unique maximizer $n^*$, strictly
decreasing in $\delta/B$. Validity at the optimum therefore rises as contests become safer, from at most $R(1)$ for a single
decisive fight toward the ceiling $\lambda^2V_P/(\lambda^2V_P+\tau^2+\rho\sigma^2)$. Equivalently, at fixed cumulative risk
$D=n\delta$ validity is $\lambda^2V_P/(\lambda^2V_P+\tau^2+\sigma^2[\rho+(1-\rho)\delta/D])$, decreasing in the severity
$\delta$.
\end{proposition}

This is the sound core of claim N2. The advantage of besting a rival ``without long-term consequence'' is not that it takes
intelligence. It is that the contest can be held again next month, against someone else, in front of the same people. A
fight to the death samples ability once and removes the loser: in the terms of Proposition~\ref{prop:identity} it is an
opportunity effect with accuracy $R(1)$. A season of wrestling is an information effect. Weapons sharpen the contrast: once
any man can kill any other from ambush or with allies, unregulated fighting becomes expensive for the strong as well as the
weak~\cite{wrangham2018,gintis2015}, $\delta$ for escalated contests rises, and regulated ones take their place.

\subsubsection{Leveling} Let a dominant male decide how many subordinate males to tolerate. Each of $k$ tolerated
subordinates fights for the group and takes a share $q$ of its reproduction, the price of his staying and fighting. The
group meets rival coalitions of strength $M$ and prevails with probability
\begin{equation}
p(k)=\frac{(1+k)^\theta}{(1+k)^\theta+M^\theta},\label{eq:csf}
\end{equation}
a ratio contest success function~\cite{hirshleifer1989} in which $\theta=1$ corresponds to Lanchester's linear law (a series
of duels) and $\theta=2$ to the square law (concentrated attack, as in raids and with
projectiles)~\cite{johnson2015,wilson2002}. The dominant's expected reproduction is $U(k)=(1-kq)\,p(k)$.

\begin{proposition}[The despot's dilemma]\label{prop:despot}
$\ln U$ is strictly concave in $k$, with a unique optimum $k^*$. (i)~$k^*>0$ if and only if\/
$\theta M^\theta/(1+M^\theta)>q$: a dominant facing weak rivals tolerates no one. (ii)~$k^*$ is increasing in the rivals'
strength $M$ and decreasing in the subordinates' price $q$. (iii)~As $M\to\infty$, $k^*\to(\theta-q)/[q(1+\theta)]$ and the
dominant's share falls to $(1+q)/(1+\theta)$: about one half under the linear law and one third under the square law.
\end{proposition}

Proposition~\ref{prop:despot} is claim N1 with its comparative statics corrected. What makes subordinates worth keeping is
not the size of one's own group but the strength of the coalitions it must face; an arms race in coalition size follows,
since each group's tolerance raises its rivals' $M$. The square law, which chimpanzee and human raiding
approximate~\cite{wilson2002,wilson2014,johnson2015}, makes the effect stronger. The model is a top-down bound in the tradition of
concession models of reproductive skew~\cite{vehrencamp1983,johnstone2000}: it gives the skew a dominant would \emph{choose}. Leveling from below, by coalitions of subordinates armed with weapons and
words~\cite{boehm1999,gavrilets2012,pandit2003}, reaches the same place faster, and the rewards that societies attach to
defense~\cite{glowacki2013} are the subordinates' price $q$ paid in standing.

\subsubsection{Synthesis: leveling re-aims selection, and slows it} By Proposition~\ref{prop:identity}, lower skew means
lower intensity. In the example of Table~\ref{tab:modelv}, the opportunity for selection among twenty males falls from
\DesIdespot{} under despotism to \DesILin{} or \DesISq{} when rivals are strong, a \DesIntLo- to \DesIntHi-fold drop in intensity.
Leveling is therefore not an accelerator. Its significance for the brain is that it changes the criterion. Where one male
can no longer take reproduction by force, standing has to be conferred by others, and others confer it on what is useful to
them: skill, judgment, generosity, the ability to rally a coalition. Under the illustrative loadings of
Table~\ref{tab:modelv}, moving standing from a contest of force to skill and prestige raises selection on capacity
\AimGain-fold and lowers selection on formidability \AimLoss-fold at equal intensity. Combined with the fall in intensity, and assuming
unchanged alignment, the net effect is a \NetGLo- to \NetGHi-fold increase in selection on capacity and a \NetFLo- to \NetFHi-fold decrease in selection
on formidability. The numbers are illustrations, but the direction is at least consistent with the hominin record, in which canine dimorphism
is minimal and body-size dimorphism modest~\cite{plavcan2012} while brains enlarged. The hero of claim N5 fits here too. A group can pay for dangerous public service
only in standing, and what it rewards is courage that is useful, which is why heroic risk attracts and gratuitous risk does
not~\cite{farthing2005}, and why the fitness of warriors varies with what their societies
reward~\cite{glowacki2013,glowacki2015,beckerman2009}.

\subsection{Model VI: episodes of selection}\label{sec:episodes}

The models above treat reproduction as one quantity. It is a sequence, and the place in the sequence where selection acts
decides who exerts it and what information they need.

\begin{proposition}[Access and number]\label{prop:episodes}
Let $c\ge0$ be completed family size, $p=\Pr(c>0)$, and relative fitness $w=c/\E c$. For any trait $z$ with mean zero,
\begin{equation}
S=\Cov(z,w)=\underbrace{\E[z\mid c>0]}_{\text{access}}+\underbrace{\frac{\Cov(z,c\mid c>0)}{\E[c\mid c>0]}}_{\text{number}},\label{eq:episodes}
\end{equation}
and the opportunity for selection splits as
\begin{equation}
I=\Var(w)=\frac{1-p}{p}+\frac{I_{c>0}}{p},\label{eq:oppsplit}
\end{equation}
where $I_{c>0}$ is the squared coefficient of variation of family size among parents.
\end{proposition}

This is the two-episode case of the partition of Arnold and Wade~\cite{arnold1984}; under natural fertility a third term, the
survival of offspring, is at least as large (Section~\ref{sec:decomp}). The two terms of Eq.~\eqref{eq:episodes} are governed by
different decisions. Access is decided mostly by other people: whether anyone accepts a person as a partner. Number is decided
mostly by the couple and their circumstances.

\begin{corollary}[A floor is enough]\label{cor:floor}
If the fraction $q$ with the lowest assessed scores do not reproduce and all others reproduce equally, then with $z_q=\Phi^{-1}(q)$
\begin{equation}
S_A=\sigma_A\,h\sqrt R\;\frac{\varphi(z_q)}{1-q},\qquad I=\frac{q}{1-q}.\label{eq:floor}
\end{equation}
The intensity $\varphi(z_q)/(1-q)$ is 0.20, 0.35 and 0.50 for $q=0.1$, 0.2 and 0.3.
\end{corollary}

A floor that excludes one or two people in ten, operated by assessors whose information is poor, delivers selection
differentials on capacity far above what the fossil rate requires (Section~\ref{sec:strength}). No champion needs to
out-reproduce anyone.

\begin{remark}[Daily life is made of easy items]
Whether a floor or a ladder can be enforced depends on what assessors can see. A task that most adults complete (bringing food
home most weeks, keeping a shelter, holding one's temper, keeping an agreement) is, in the language of test theory, an easy
item: its outcome carries information about people near the bottom of the distribution and almost none about differences among
the able. Hard items that separate the best from the good are rare events with few observers. The information that
accumulates about everyone, at no cost, is therefore information about the lower tail, which is the information a floor needs.
\end{remark}

\subsection{Model VII: capacity and delivery}\label{sec:delivery}

A partner's value to a chooser is not the partner's capacity alone. Let $b$ be the benefit a candidate can produce, and
$p\in[0,1]$ the share of it that reaches the chooser and the chooser's children: willingness and reliability of investment.
With compatibility $c$ and expected harm $h$, a chooser's expected return is $v=pb+c-h$. The division follows the factors found
in studies of partner ideals~\cite{fletcher1999} and the logic of choice for direct benefits~\cite{hoelzer1989}. Write
$b=\bar b+B$ and $p=\bar p+P$ with $B$ and $P$ heritable, independent and jointly normal with their breeding values, and let
reproductive contribution rise with assessed value, $w=1+s(\hat v-\E\hat v)$, where capacity and delivery are assessed with
validities $R_b$ and $R_p$.

\begin{proposition}[Capacity pays in proportion to delivery]\label{prop:delivery}
The selection differentials on the two breeding values are
\begin{equation}
S_{A_B}=s\,\bar p\,R_b\,V_{A,B},\qquad S_{A_P}=s\,\bar b\,R_p\,V_{A,P}.\label{eq:delivery}
\end{equation}
With quadratic costs $\tfrac12c_B\bar b^2$ and $\tfrac12c_P\bar p^2$, the state without investment or surplus capacity
($\bar b=\bar p=0$) is stable if $s^2R_bR_p<c_Bc_P$ and unstable otherwise; when it is unstable the two traits rise together
until delivery reaches its bound, where $\bar b^*=sR_b/c_B$.
\end{proposition}

Three consequences follow. First, where nothing a partner produces reaches the chooser ($\bar p\approx0$), choice for
benefits exerts no selection on capacity, however able the candidates and however accurately ability is seen. Capacity
becomes a target of choice only in a species in which partners invest, and human paternal investment is unusual among
mammals~\cite{geary2000,marlowe2000}. Second, the two traits are complements: each strengthens selection on the other, so
capacity and commitment take off together or not at all. Third, they are known through different evidence. Capacity shows in
performance, which is public and quick to observe. Delivery shows in conduct when interests conflict, which is private and
slow. $R_p$ is raised by long acquaintance and by assessors with a stake, such as parents and kin, and not by spectacle. As
history the joint take-off is conjecture; the complementarity is testable now (prediction P23).

\subsection{Model VIII: tests, enforcement and verification}\label{sec:verification}

Models I--VII take validity as given. A candidate who knows when assessment is taking place can change what it shows, and
the two components of Model VII respond in opposite ways. Capacity cannot be displayed beyond its level, so the occasions on
which a candidate tries hardest are the best evidence of it. Delivery is a choice, so an occasion on which the candidate
knows the choice is being judged is no evidence of what will be chosen later. Students of animal signals draw the same line
between signals that cannot be faked and signals that must be kept honest by cost or checking~\cite{maynardsmith1995,grafen1990}.
Occupational psychology separates maximum from typical performance, which correlate only moderately, with ability the
stronger correlate of the first~\cite{sackett1988,beus2012}. And people reason about one another this way: perceivers ask
first whether another person intends good or ill and then whether that person is able to act on the
intention~\cite{fiske2007}, and they weigh good performances more heavily when judging competence and bad conduct more
heavily when judging morality~\cite{reeder1979,skowronski1987,wojciszke1993}.

Let an assessor observe $n$ occasions, of which a share $d$ are \emph{staged}, that is, recognizable to the candidate as
assessment. On a staged occasion the candidate performs at capacity, $y=b+\varepsilon$ with $\Var\varepsilon=\sigma^2$, and
delivers. On an unstaged occasion the candidate defaults with probability $q=1-p$, where $p$ is the disposition of
Model~VII with variance $V_P$; write $\nu=\E[q(1-q)]$. The assessor scores capacity by the mean of the staged performances
and delivery by the number of defaults seen.

\begin{proposition}[A staged test measures capacity; unstaged observation measures delivery]\label{prop:staging}
(i)~The validities for capacity and for delivery are
\begin{equation}
R_b=\frac{V_B}{V_B+\sigma^2/(dn)},\qquad R_p=\frac{V_P}{V_P+\nu/[(1-d)n]},\label{eq:staging}
\end{equation}
so $R_b$ rises and $R_p$ falls with $d$, and $R_p\to0$ as $d\to1$ for every $n$. (ii)~A candidate with default rate $q$
keeps a clean record with probability $(1-q)^{(1-d)n}$. An even chance of seeing one default needs
$(1-d)n\ge\ln2/[-\ln(1-q)]$ occasions, and after a clean record the 95\% upper limit on $q$ is
$1-0.05^{1/[(1-d)n]}\approx3/[(1-d)n]$. (iii)~The product $R_bR_p$, which sets the take-off condition of
Proposition~\ref{prop:delivery}, vanishes at $d=0$ and at $d=1$ and is greatest at
\begin{equation}
d^*=\frac{\sqrt{A(1+C)}}{\sqrt{A(1+C)}+\sqrt{C(1+A)}},\qquad A=\frac{\sigma^2}{nV_B},\quad C=\frac{\nu}{nV_P}.\label{eq:dstar}
\end{equation}
\end{proposition}

Only occasions that the candidate cannot tell from ordinary life carry information about delivery, and no number of
recognized tests substitutes for them. Part~(ii) is the rule of three for a record without failures~\cite{hanley1983}:
reliability is never established, only not yet refuted, at a rate set by unstaged observation alone. The result is the
validity form of an audit-volume bound derived for the oversight of artificial agents in a companion
manuscript~\cite{taylor2026iina}, where the same asymmetry between what a system can do and what it will do appears. With
the illustrative values of experiment E16 (\VerN{} occasions), moving from one staged occasion in ten to nine in ten raises
validity for capacity from \VerRbLo{} to \VerRbHi{} and lowers validity for delivery from \VerRpLo{} to \VerRpHi{}
(Fig.~\ref{fig:main}D). A candidate who defaults on a quarter of unstaged occasions then keeps a clean record with
probability \VerCleanNine\%, against \VerCleanHalf\% when half of the occasions are unstaged. Part~(iii) rests on the
simplification that capacity is read only from staged occasions; if daily life also shows capacity, $d^*$ is smaller and
parts (i) and (ii) are unchanged for delivery. Either way an assessor needs both kinds of occasion: the contest or festival
for what a person can do, and unremarkable days for what the person does.

Three familiar practices raise $(1-d)n$. One is time: courtship has been modeled as time traded for information about a
partner's intentions~\cite{wachtmeister1999,seymour2009,sozou2005}, about two-thirds of relationships begin as
friendships~\cite{stinson2022}, and brideservice keeps a suitor under the eyes of his in-laws before a marriage is secure
(Section~\ref{sec:evidence}). The second is scale: in a band of a few dozen, conduct is seen daily by people who are not
its audience. The third is report.

\begin{remark}[What is shown cannot reveal what is withheld]
A candidate also chooses which occasions an assessor attends. Any statistic computed from the attended occasions has the
same distribution for a candidate with nothing to hide as for one who keeps the same faults away from every attended
occasion, so no amount of attended observation bounds what was concealed. The remedy is an informant whose interests differ
from the candidate's: a rival, who gains by reporting faults~\cite{buss1990}; a former partner; the chooser's own kin. If
each of $m$ informants independently knows of a concealed fault with probability $\rho$, it stays hidden with probability
$(1-\rho)^m$; informants who share one source count as one (the shared distortion $\varsigma^2$ of
Section~\ref{sec:portfolio}). Deception is a documented tactic in courtship~\cite{tooke1991}, women tend to underestimate
men's commitment~\cite{haselton2000}, and moral character outweighs warmth in the impressions people
form~\cite{goodwin2014}: assessors behave as if they know that display is selected.
\end{remark}

\subsubsection{Enforcement} Verification is one answer to a disposition that cannot be read. The other is to make the
disposition matter less. Let a share $e\in[0,1]$ of every candidate's product reach the assessors whatever the candidate's
disposition, because third parties enforce it. Among foragers the meat of large animals is widely shared and much sharing is
by demand~\cite{woodburn1982,peterson1993,gurven2004}; norms are enforced by criticism, ridicule and
exclusion~\cite{wiessner2005,boehm1999}; and a successful Hadza hunter has been reported not to control the distribution
of his kill~\cite{hawkes2001}. The delivered share is then $\delta=e+(1-e)p$ and a candidate's value is $\delta b$.

\begin{proposition}[Enforcement substitutes for verification]\label{prop:enforce}
(i)~The selection differentials are
\begin{equation}
S_{A_B}=s\,[e+(1-e)\bar p]\,R_b\,V_{A,B},\qquad S_{A_P}=s\,(1-e)\,\bar b\,R_p\,V_{A,P}.\label{eq:enforce}
\end{equation}
(ii)~With the costs of Proposition~\ref{prop:delivery} and $T=s^2R_bR_p/(c_Bc_P)$, the state $\bar b=\bar p=0$ is not an
equilibrium for any $e>0$. If $T(1-e)<1$ the unique stable equilibrium is
\begin{equation}
\bar b^*=\frac{sR_b}{c_B}\cdot\frac{e}{1-T(1-e)^2},\qquad \bar p^*=\frac{T\,e\,(1-e)}{1-T(1-e)^2};\label{eq:enforceeq}
\end{equation}
otherwise delivery reaches its bound and $\bar b^*=sR_b/c_B$. (iii)~For $T<1$, $\bar b^*$ increases with $e$ from zero to
$sR_b/c_B$, while $\bar p^*$ is zero at $e=0$ and at $e=1$ and greatest at $e^*=1-(1-\sqrt{1-T})/T$.
\end{proposition}

Three consequences follow. First, enforcement removes the threshold. Proposition~\ref{prop:delivery} left a circle:
capacity is worth choosing only where partners deliver and delivery only where partners are capable, so below $T=1$ nothing
starts. With $e>0$ capacity is selected from the outset, in proportion to what is enforced and with no evolved commitment
at all. In the illustration of E16, with $T=\EnfT$, enforcing one-fifth of the product brings capacity to \EnfCapFifth{} of
its maximum and enforcing half brings it to \EnfCapHalf{} (Fig.~\ref{fig:verify}A). Second, enforcement hides disposition
from selection: $S_{A_P}$ carries the factor $(1-e)$. Evolved commitment is therefore greatest under partial enforcement
($e^*=\EnfEstar$ in the illustration) and is not selected at all through choice when enforcement is complete; persistent
takers can still be selected against directly, by the sanctions themselves~\cite{wrangham2019}. Where choice alone is
already above threshold the same factor works against delivery: with $T=\EnfCrowdT$, enforcing \EnfCrowdAt{} of the product
lowers the delivered share at equilibrium from one to \EnfCrowdMin{} (Fig.~\ref{fig:verify}B), because what is enforced no
longer selects for commitment. Third, enforcement changes
who the assessors are. Product that is shared across a band reaches everyone, so a wife gains little by marrying the
producer, which is the puzzle that the show-off hypothesis addressed~\cite{hawkes2002}. Under enforcement the assessors
whose decisions matter are those who receive the product and control access: the people who confer standing, and the
parents and kin who arrange most marriages (Section~\ref{sec:evidence}). How much of a hunter's product his own household
keeps is disputed: later Hadza data show sharing that favors producers' households, with the best hunters providing three
to four times as much food to their families as median hunters~\cite{wood2013,gurvenhill2009}. In the model that dispute
concerns the size of $e$ in one population, not the logic.

The order of events that this suggests, limits on what the strong could keep first and selection on capacity afterward, is
conjecture, as is the leveling synthesis of Section~\ref{sec:modelv}, which says the same of the criterion of
standing. The comparative statics can be tested now (predictions P24--P26 and P28).

\begin{corollary}[Breadth taxes exaggeration]\label{cor:inflation}
Let candidates be able to raise the appearance of performance in domain $k$ by $f_k$, at cost $\tfrac12c_Ff_k^2$ in log
fitness and with no benefit to anyone. In the setting of Proposition~\ref{prop:polymath}, exaggeration is selected as a
specialty is: its equilibrium level is $f_k^*=\beta_S/(K\,c_F\,\SD(Y))$ per domain, the total spent on it is proportional
to $1/K$, and heritable variation in domain-specific showmanship is selected, at equal genetic variance, $K$ times more
weakly than general capacity (Eq.~\ref{eq:aimK}). A talent for seeming that operates in every domain is not diluted; it belongs to the general nuisance
$\Omega$ of Eq.~\eqref{eq:RG}.
\end{corollary}

Breadth therefore protects an assessment against gaming without requiring that display be costly: a candidate who
inflates one performance moves a $K$-domain average by $1/K$ of the inflation, while real general capacity moves it in
full. What breadth cannot remove is an impression that carries across every venue. That takes evidence of a different
kind, the record of Section~\ref{sec:decomp}: what was produced and who received it, kept by people with a stake in the
answer. An indicator on which decisions depend comes under pressure to be
gamed~\cite{campbell1979}; common comparison across many domains, records kept by interested parties, and enforcement by
third parties are three defenses that do not require reading anyone's motives.

\subsection{Model IX: who decides, who becomes an ancestor, and how long a choice must wait}\label{sec:modelix}

Four further results answer questions raised in review of the previous version (Section~\ref{sec:notes}). Each is short,
and each was checked by simulation (E21, Section~\ref{sec:experiments}).

\subsubsection{Coercion} Let a share $\kappa$ of matings be settled without regard to a trait $A$ that assessors would
otherwise use, and let the rest be settled by assessment. Fitness is a mixture, so the differential on $A$ is the mixture
of the differentials, Eq.~\eqref{eq:coercion}. If what settles the coerced matings is uncorrelated with $A$, then
$S_A=(1-\kappa)S_{\text{assessed}}$ exactly: in simulation, with the upper half of a population reproducing and
$\kappa=0$, 0.25, 0.5, 0.75, the differential on the assessed trait falls in proportion and that on the trait rewarded by
force rises by the same amount (largest error \NtCoeErr). The result needs no assumption about who assesses. It applies as
well when parents choose, with $S_{\text{assessed}}$ computed for the trait their criterion loads on.

\subsubsection{Effective numbers}
\begin{proposition}[Who becomes an ancestor]\label{prop:ancestors}
In a large population, let a share $q$ of one sex leave no offspring, and let the number left by the others have mean $m$
and variance $m+\alpha m^2$, where $\alpha=0$ is the Poisson case and $\alpha>0$ measures lasting differences among
parents. Then the variance effective number of that sex, for an autosomal or a uniparentally inherited marker, is
\begin{equation*}
\frac{N_e}{N}=\frac{\bar k}{\bar k-1+V/\bar k}=\frac{1-q}{1+\alpha},
\end{equation*}
where $\bar k$ and $V$ are the mean and variance over all members of the sex. The ratio does not depend on mean family
size. For two sexes of equal number,
$N_{e,f}/N_{e,m}=\big[(1-q_f)/(1-q_m)\big]\big[(1+\alpha_m)/(1+\alpha_f)\big]$.
\end{proposition}

A ratio of two is therefore not a count. With $q_f=0.2$ it is met by $q_m=0.6$ and equal $\alpha$ (the popular ``80\% of
women, 40\% of men''); by $q_m=0.2$ with $\alpha_m=1$ and $\alpha_f=0$; and by every combination between
(Fig.~\ref{fig:choice}D). It is also met with no difference in individual reproduction at all, when men compete as members
of patrilineal groups that grow or die out together, since the lineage and not the man is then the unit that
drifts~\cite{zeng2018}; and it is moved by sex differences in migration and generation length, which the proposition
ignores. In the General Social Survey (completed families, \VgAncN{} respondents born 1900--1979) the directly computed
ratio and the identity agree to rounding error ($\AudAnc$), $\alpha$ is slightly negative in both sexes (\VgAlphaM{} for men
and \VgAlphaW{} for women: families are more equal than chance), and $N_{e,f}/N_{e,m}=\VgNeRatio$ (\VgNeRatioCI). The
standardized variance in number of children, Crow's index~\cite{crow1958}, is \VgIm{} for men and \VgIf{} for women. Men report somewhat
fewer children than women do in surveys, which biases $q_m$ upward.

\subsubsection{Waiting for evidence}
\begin{proposition}[Thin evidence delays commitment]\label{prop:waiting}
Suppose a commitment is made when a partner's delivery is known with reliability $R^*$. Let each of $k$ independent
witnesses observe $\lambda$ occasions per period, of which a share $d$ are recognized by the candidate as tests and carry
no information about delivery (Proposition~\ref{prop:staging}), and let one unstaged observation have reliability $r_1$.
Then the waiting time is
\begin{equation}
T=\frac{R^*(1-r_1)}{r_1(1-R^*)}\cdot\frac{1}{k\,\lambda\,(1-d)}.\label{eq:waiting}
\end{equation}
It is inversely proportional to the number of witnesses and to the unstaged share. If births begin at commitment and are
spaced $b$ periods apart within a fertile window $W$ that opens at age $a_0$, completed family size is
$\max\{0,(W-a_0-T)/b\}$, which falls linearly in $T$ and reaches zero at $T=W-a_0$.
\end{proposition}

With $r_1=0.05$, $R^*=0.8$ and four occasions per witness per period, a camp of thirty in which a fifth of occasions are
staged reaches the threshold in \NtWaitCamp{} periods, five friends in \NtWaitFive, one observer in \NtWaitOne, and one
observer who sees mostly staged occasions in \NtWaitOneStaged{} (simulated reliabilities within \NtWaitErr{} of $R^*$).
The numbers are illustrations. The proportionality is the result, and it is what connects a thinner community to later
and fewer unions without any change in what people want. It is a conjecture about present-day populations. One finding
runs against its simplest form, that couples who met online in the United States married sooner~\cite{rosenfeld2017}; those
who search online may be selected for readiness to commit, and the proposition is about the wait before a couple forms as
much as after.

\subsubsection{Safe contests and narrow tests}
\begin{corollary}[Safe contests buy reliability by repetition]\label{cor:ritual}
Let one dangerous trial reveal capacity with reliability $\rho_d$ and one safe trial with reliability $\rho_s<\rho_d$,
errors independent across trials. Then $n$ safe trials are at least as reliable about capacity as the dangerous one if and
only if $n\ge\rho_d(1-\rho_s)/[\rho_s(1-\rho_d)]$. If the safe trials are recognized as tests, their validity for
delivery under risk is zero for every $n$ (Proposition~\ref{prop:staging}).
\end{corollary}

\begin{corollary}[A narrow test makes a deficit of a difference]\label{cor:narrow}
Let a trait $x$ load $a_k$ on performance in domain $k$, and let assessment use the average of a set $T$ of domains. The
selection, or evaluation, gradient on $x$ is proportional to $\bar a_T$, its mean loading over $T$. A trait with
$\sum_{k=1}^{K}a_k=0$ is neutral under comparison on all $K$ domains and is penalized in proportion to $|\bar a_T|$ under a
test confined to domains on which it hurts, while the gradient on general ability is the same under both.
\end{corollary}

In simulation with ten domains, a trait that adds 0.3 standard deviations in five and subtracts 0.3 in the other five has
a slope of \NtNarBroad{} on the ten-domain average and \NtNarHurt{} on the average of the five on which it hurts; the slope
on general ability is one in both. The corollary is linear algebra, and it states exactly what the generalist result
implies for variation that is not general: breadth of comparison is what keeps mixed profiles neutral, and narrowing the
comparison is what turns one of them into a deficit.
